\section{Purpose}

Every failure this project has hit shared one shape: a formula that is correct
in the common case broke silently at an edge case, and the break wasn't caught
until a human recomputed one number by hand and noticed it didn't match. This
document exists to make that hand-check repeatable: three physical quantities,
their formulas, why each formula is what it is, and one fully worked, verified
example you can reproduce with a calculator.

\section{The Three Quantities That Decide a Winner}

Every candidate array design is judged on exactly three physical properties.
All three depend only on the hole coordinates, the hole radius, and the pupil
radius --- nothing exotic.

\begin{enumerate}
  \item \textbf{Viability} --- does any light get through at all, at the
  gaze angle being checked?
  \item \textbf{Ghost margin} --- if two holes are both illuminated at once,
  do they produce a double image?
  \item \textbf{Packing density} --- how many holes fit per unit area, i.e.\
  how much of the design actually contributes light?
\end{enumerate}

A design only counts as a real answer if all three check out --- and the
entire point of a hand-check is that each of these three is simple enough to
compute on paper for one design in a couple of minutes.

\newpage
